\documentclass[11pt]{amsart}

\usepackage[T1]{fontenc}
\usepackage{lmodern}
\usepackage{microtype}
\usepackage{amsmath,amssymb,amsthm}
\usepackage[hidelinks]{hyperref}

\hypersetup{
  pdftitle={Counterexamples to Wirdane's Cyclotomic Catalan Conjecture}
}

\newtheorem{theorem}{Theorem}
\theoremstyle{remark}
\newtheorem*{remark}{Remark}

\newcommand{\Cat}{\operatorname{Cat}}

\title[Counterexamples to a Cyclotomic Catalan Conjecture]
{Counterexamples to Wirdane's Cyclotomic Catalan Conjecture}

\subjclass[2020]{Primary 05A15; Secondary 05A30}
\keywords{$q$-Catalan numbers, roots of unity, cyclotomic specialization}
\date{}

\begin{document}

\begin{abstract}
We disprove the literal form of Wirdane's Cyclotomic Catalan
Conjecture, Conjecture~6.6 in arXiv:2605.14682v1. For nontrivial
cyclotomic orders, the first counterexample in the parameter $n$ is
$(n,\mu)=(2,3)$. We also give three infinite families of
counterexamples: $(2m+1,2)$ for $m\ge1$, $(3r+2,3)$ for $r\ge0$,
and $(p-1,p)$ for every odd prime $p$.
\end{abstract}

\maketitle

Let $C_0(q)=1$ and, for $n\ge1$, let
\begin{equation}\label{eq:recurrence}
  C_n(q)=\sum_{k=0}^{n-1}q^k C_k(q)C_{n-1-k}(q)
\end{equation}
be the F\"urlinger--Hofbauer $q$-Catalan numbers~\cite{FH}. Write
\[
  \Cat_m=\frac{1}{m+1}\binom{2m}{m},
  \qquad
  \mathcal C(t)=\sum_{m\ge0}\Cat_m t^m
               =1+t\mathcal C(t)^2.
\]
At $q=1$, recurrence~\eqref{eq:recurrence} is the Catalan recurrence,
so $C_n(1)=\Cat_n$.

Fix $\mu\ge2$ and $\zeta_\mu=e^{2\pi i/\mu}$. We first recall the
objects in terms of which Conjecture~6.6 of~\cite{Wirdane} is stated.
Let $S_n(312)$ be the set of $312$-avoiding permutations of
$\{1,\ldots,n\}$, let $S_{n,k}(312)=\{\pi\in S_n(312):\pi_{k+1}=n\}$
for $0\le k\le n$, and let
$S'_{n,k}(312)=\bigsqcup_{j=0}^{k}S_{n,j}(312)$. The value $n$ occupies
exactly one position of each $\pi\in S_n(312)$, so
$S_{n,0}(312),\ldots,S_{n,n-1}(312)$ partition $S_n(312)$ while
$S_{n,n}(312)=\varnothing$; hence
\begin{equation}\label{eq:sprime}
  S'_{n,n}(312)=S_n(312).
\end{equation}
For $1\le i\le\mu$, the $i$-inversions of $\pi$ are counted by
\[
  \operatorname{inv}_i(\pi)
  =\#\bigl\{(a,b):1\le a<b\le n,\ \pi_a>\pi_b,\
      \pi_a\equiv i\ (\mathrm{mod}\ \mu)\bigr\}
\]
\cite[Definition~6.1]{Wirdane}, and the multivariate $q$-Catalan
triangle is
\[
  C_{n,k}(q_1,\ldots,q_\mu)
  =\sum_{\pi\in S'_{n,k}(312)}\prod_{i=1}^{\mu}
     q_i^{\operatorname{inv}_i(\pi)}
\]
\cite[Definition~6.2]{Wirdane}, its cyclotomic analogue being
$C_{n,k}^{(\mu)}=C_{n,k}(\zeta_\mu,\ldots,\zeta_\mu)$. Because
$1,\ldots,\mu$ is a complete residue system modulo $\mu$, each
inversion $(a,b)$ of $\pi$ is counted by exactly one of the
$\operatorname{inv}_i$, so
$\sum_{i=1}^{\mu}\operatorname{inv}_i(\pi)=\operatorname{inv}(\pi)$,
the inversion number. With~\eqref{eq:sprime} this gives
\begin{align}\label{eq:diagonal}
  C_{n,n}^{(\mu)}
  &=\sum_{\pi\in S_n(312)}
       \zeta_\mu^{\sum_{i=1}^{\mu}\operatorname{inv}_i(\pi)}
    =\sum_{\pi\in S_n(312)}
       \zeta_\mu^{\operatorname{inv}(\pi)}
    =C_n(\zeta_\mu),
\end{align}
the last equality being the interpretation
\[
  C_n(q)=\sum_{\pi\in S_n(312)}q^{\operatorname{inv}(\pi)}
\]
of~\cite{FH}, recorded in~\cite{Wirdane} as equation~(2) and again as
Remark~3.4. Equivalently, one may specialize all $q_i$ to a common $q$
by \cite[Corollary~6.5]{Wirdane} and then invoke
\cite[Theorem~3.3]{Wirdane}, namely
$C_{n,k}(q)=\sum_{\pi\in S'_{n,k}(312)}q^{\operatorname{inv}(\pi)}$.
We do not use \cite[Proposition~2.4]{Wirdane}, which asserts
$C_{n,n}(q)=C_n(q)$ outright, its proof being deferred there to a work
in preparation.
Conjecture~6.6 of~\cite{Wirdane} reads, in full: ``If $\mu\mid(n+1)$:
$C_{n,n}^{(\mu)}=\frac{1}{\mu}\binom{n+1}{(n+1)/\mu}$.'' By
\eqref{eq:diagonal} it asserts that, whenever $\mu\mid n+1$,
\begin{equation}\label{eq:conjecture}
  C_n(\zeta_\mu)
  =\frac{1}{\mu}\binom{n+1}{(n+1)/\mu}.
\end{equation}
As printed the conjecture places no restriction on $\mu$, and the case
$\mu=1$ is discussed in \cite[Remark~6.4]{Wirdane}; we standardize on
$\mu\ge2$ because $\mu=1$ is a degenerate reading of ``cyclotomic'',
and dispose of $\mu=1$ separately in the closing Remark.

\begin{theorem}\label{thm:main}
Equation~\eqref{eq:conjecture} is false. Each of the following is an
infinite family of counterexamples:
\[
  (n,\mu)=(2m+1,2)\quad(m\ge1),
  \qquad
  (n,\mu)=(3r+2,3)\quad(r\ge0),
\]
\[
  (n,\mu)=(p-1,p)\quad(p\text{ an odd prime}).
\]
Among $\mu\ge2$, the least value of $n$ occurring in a counterexample
is $n=2$, with $\mu=3$.
\end{theorem}

\begin{proof}
All generating functions below are formal power series. Set
\[
  F_q(z)=\sum_{n\ge0}C_n(q)z^n.
\]
Recurrence~\eqref{eq:recurrence} is equivalent to
\begin{equation}\label{eq:functional}
  F_q(z)=1+zF_q(z)F_q(qz).
\end{equation}

\medskip
\noindent\emph{Order two.}
Put $F(z)=F_{-1}(z)$. Applying~\eqref{eq:functional} at $z$ and
$-z$ gives
\[
  F(z)=1+zF(z)F(-z),
  \qquad
  F(-z)=1-zF(-z)F(z),
\]
so $F(z)+F(-z)=2$. Thus, for $H(z)=F(z)-1$,
\[
  H(z)=z\bigl(1-H(z)^2\bigr).
\]
The unique solution with zero constant term is
$H(z)=z\mathcal C(-z^2)$. Consequently,
\begin{equation}\label{eq:minusone}
  C_{2m}(-1)=0\quad(m\ge1),
  \qquad
  C_{2m+1}(-1)=(-1)^m\Cat_m\quad(m\ge0).
\end{equation}
For $n=2m+1$, the conjectured value is
\[
  \frac12\binom{2m+2}{m+1}=(2m+1)\Cat_m,
\]
which differs from~\eqref{eq:minusone} for every $m\ge1$.

\medskip
\noindent\emph{Order three.}
Let $\omega=\zeta_3$ and put $A(z)=\mathcal C(z^3)$. Define
\[
  G(z)=\frac{1+\omega z^2A(z)}{1-z}.
\]
Since $A(z)=1+z^3A(z)^2$ and $1+\omega+\omega^2=0$,
\[
  (1+\omega zA(z))(1-\omega z)
  =(1+\omega z^2A(z))(1+z^2A(z)).
\]
Moreover,
\[
  G(\omega z)=\frac{1+z^2A(z)}{1-\omega z},
\]
so the displayed identity gives $G(z)-1=zG(z)G(\omega z)$. The
functional equation~\eqref{eq:functional} determines its coefficients
recursively, so
\begin{equation}\label{eq:omegaGF}
  F_\omega(z)=\frac{1+\omega z^2\mathcal C(z^3)}{1-z}.
\end{equation}
Therefore, for $r\ge0$,
\begin{equation}\label{eq:omegaeval}
  C_{3r+2}(\omega)
  =1+\omega\sum_{j=0}^{r}\Cat_j.
\end{equation}
Its imaginary part is positive, whereas the right side of
\eqref{eq:conjecture} is real. Thus every $(3r+2,3)$ is a
counterexample. In particular,
\[
  C_2(\omega)=1+\omega\ne1
  =\frac13\binom31.
\]

\medskip
\noindent\emph{Odd prime orders.}
Let $p$ be an odd prime, $\zeta=\zeta_p$, and
$\mathfrak p=(1-\zeta)\subset\mathbb Z[\zeta]$. Since
$\Phi_p(1)=p$,
\[
  \mathbb Z[\zeta]/\mathfrak p\cong\mathbb F_p,
  \qquad \zeta\equiv1\pmod{\mathfrak p}.
\]
Hence
\[
  C_{p-1}(\zeta)
  \equiv C_{p-1}(1)
  =\Cat_{p-1}\pmod{\mathfrak p}.
\]
Moreover,
\[
  (2p-1)\Cat_{p-1}
  =\binom{2p-1}{p-1}
  =\prod_{j=1}^{p-1}\frac{p+j}{j}
  \equiv1\pmod p,
\]
so $\Cat_{p-1}\equiv-1\pmod p$. Therefore
\[
  C_{p-1}(\zeta_p)\equiv-1\not\equiv1\pmod{\mathfrak p},
\]
the two residues being distinct because $p$ is odd. The value
predicted by~\eqref{eq:conjecture} at $(n,\mu)=(p-1,p)$
is $\frac1p\binom p1=1$, proving the prime family.

Finally, no $\mu\ge2$ divides $n+1$ when $n=0$. When $n=1$, the
only admissible order is $\mu=2$, and both sides of
\eqref{eq:conjecture} equal $1$. Thus $(2,3)$ has the least possible
value of $n$.
\end{proof}

\begin{remark}
The three families are far from exhausting the failures. Among the
$81$ pairs $(n,\mu)$ with $\mu\ge2$, $\mu\mid n+1$ and $n\le29$,
reducing $C_n(q)$ modulo the cyclotomic polynomial $\Phi_\mu(q)$
shows that~\eqref{eq:conjecture} holds at exactly one of them,
$(n,\mu)=(1,2)$, and fails at the remaining $80$; the three families
of Theorem~\ref{thm:main} account for $32$ of the $81$. The
order-three argument is also the cheapest of the three, using only
that the left side of~\eqref{eq:conjecture} is not real while the
right side is; that observation alone disposes of $66$ of the $81$
pairs, against the $9$ reached by the congruence for odd prime orders.
\end{remark}

\begin{remark}
For $\mu=2$, Conjecture~6.6 is in conflict with an assertion made one
line later in~\cite{Wirdane}. Remark~6.7 there states: ``For $\mu=2$:
$C_{n,k}^{(2)}=(-1)^{\lfloor n/2\rfloor}C_{n,k}$ (verified for
$n\le5$).'' Here $C_{n,k}$ is the classical Catalan triangle, obtained
by setting every $q_i=1$ in Definition~6.2
\cite[Corollary~6.5]{Wirdane}; the computation
giving~\eqref{eq:diagonal} then yields $C_{n,n}=C_n(1)=\Cat_n$ on the
diagonal, so Remark~6.7 asserts that
$C_n(-1)=(-1)^{\lfloor n/2\rfloor}\Cat_n$. This is false for every
even $n\ge2$, because $\Cat_n\ne0$ whereas $C_n(-1)=0$
by~\eqref{eq:minusone}; the smallest instance is $n=2$, where it
predicts $-\Cat_2=-2$ against $C_2(-1)=0$. At $n=k=3$ three values
collide: the true one is $C_{3,3}^{(2)}=C_3(-1)=-1$, visible both
from~\eqref{eq:minusone} and from the value
$C_{3,3}(q)=1+2q+q^2+q^3$ tabulated in~\cite{Wirdane}; Remark~6.7
predicts $-\Cat_3=-5$; and Conjecture~6.6 predicts
$\frac12\binom42=3$. The $\mu=2$ branch of Conjecture~6.6 is
therefore already contradicted inside~\cite{Wirdane}, independently of
the counterexamples above.
\end{remark}

\begin{remark}
Reference~\cite{Wirdane} commits to the
Carlitz--F\"urlinger--Hofbauer normalization, through its equation~(1),
which is~\eqref{eq:recurrence}, and through its Theorem~3.3 and
Proposition~2.4. Take instead the Gaussian $q$-Catalan number
\[
  \widetilde{\Cat}_n(q)=\frac{1}{[n+1]_q}\binom{2n}{n}_{\!q},
  \qquad [j]_q=1+q+\cdots+q^{j-1},
\]
where $\binom{2n}{n}_{\!q}$ is the Gaussian binomial coefficient. Then
the right-hand side of~\eqref{eq:conjecture} is reproduced exactly at
$\mu=2$: one has
$\widetilde{\Cat}_n(-1)=\frac12\binom{n+1}{(n+1)/2}$ for every odd
$n\le11$, the common values being $1$, $3$, $10$, $35$, $126$, $462$.
At order three, by contrast,
$\widetilde{\Cat}_2(q)=1+q^2$, so
$\widetilde{\Cat}_2(\zeta_3)=1+\zeta_3^2=\frac12-\frac{\sqrt3}{2}i
\ne1$.
\end{remark}

\begin{remark}
The theorem concerns the literal formulation of Conjecture~6.6 in
arXiv:2605.14682v1, with the cyclotomic specialization derived directly
from Definitions~6.1 and~6.2 as in~\eqref{eq:diagonal}; no other
specialization assertion in~\cite{Wirdane} is used. The theorem makes
no assertion about a later revised formulation. If the trivial order $\mu=1$ is
admitted, the same formula already fails at $n=2$. The case $p=3$ in
the prime family coincides with the first member of the order-three
family.
\end{remark}

\begin{thebibliography}{9}

\bibitem{FH}
J.~F\"urlinger and J.~Hofbauer,
\emph{$q$-Catalan numbers},
J. Combin. Theory Ser. A \textbf{40} (1985), no.~2, 248--264.
\href{https://doi.org/10.1016/0097-3165(85)90089-5}
{doi:10.1016/0097-3165(85)90089-5}.

\bibitem{Wirdane}
Y.~Wirdane,
\emph{Combinatorial study of the $q$-Catalan triangle and its
generalizations},
arXiv:2605.14682v1 [math.CO], 2026.
\href{https://arxiv.org/abs/2605.14682v1}
{arXiv:2605.14682v1}.

\end{thebibliography}

\end{document}